\documentclass[11pt]{article}
\usepackage[margin=1in]{geometry}
\usepackage{amsmath,amssymb}
\usepackage{xcolor}
\usepackage[hidelinks]{hyperref}
\usepackage{enumitem}
\setlist{leftmargin=*,itemsep=0.5em}
\newcommand{\manuscript}{\textit{Construction of Kochen--Specker Sets from Mutually Unbiased Bases}}

\begin{document}

\begin{center}
{\Large Response to the Referee}\\[0.5em]
\manuscript\\[0.25em]
Mirko Navara and Karl Svozil\\[0.25em]
4 August 2026
\end{center}

Dear Referee,

Thank you for your careful and constructive report.  We agree that the principal presentation issues you identified were well founded.  We have substantially rewritten and reorganized the manuscript, and all added or changed text is shown in blue in the revised version.  We have also carried out a new exact audit of the ray and context inventories.  That audit led to several substantive corrections beyond the requested expository changes.

Our point-by-point response is as follows.

\begin{enumerate}
\item \textbf{Terminology concerning contexts and intertwining vectors.}

We now define a ray, a context (block), the associated uniform hypergraph, and an intertwining ray at the beginning of the Introduction.  A second definition states the exact-one condition for a two-valued state and distinguishes a KS configuration, a separating family of states, and unitality.  We also define maximal unbiasedness before the higher-dimensional constructions.

\item \textbf{Length and focus of the Introduction.}

The Introduction has been rewritten and shortened.  The former extended historical and polemical discussion has been condensed to a brief explanation of why the incidence structure of contexts must be specified.  The revised Introduction states the two contributions in a numbered list and ends with a concise road map.

\item \textbf{Visibility of definitions, constructions, and main results.}

The central claims are now isolated as explicit definitions and propositions.  Proposition~1 gives the exact qutrit ray and context inventory and explains the verification method.  Propositions~2 and~3 state and prove the $D=4$ and $D=5$ forcing constraints.  The distinction between a ray set, a selected context hypergraph, and the complete orthogonality hypergraph is now made explicit throughout.

\item \textbf{Guiding overview for the larger construction.}

At the start of the qutrit section we added a four-stage guide: seed selection, application of the 25 algebraic maps, projective identification and orthogonality enumeration, and complete provenance classification.  At the start of the higher-dimensional section we similarly separate the scaffold, connector constraints, and context completions.

\item \textbf{Higher-dimensional constraints versus computational detail.}

The forcing equations are now stated and proved before the coordinate inventories.  For $D=4$, the three connector equations are displayed together and their equal-modulus consequence is proved in Proposition~2.  For $D=5$, the four equations $\langle u,g_i\rangle=|x_i|^2-|x_5|^2=0$ are isolated in Proposition~3.  Completion rays are described separately from the constraints they implement.
\end{enumerate}

During the revision we also made the following author-initiated corrections.

\begin{itemize}
\item \textbf{Corrected qutrit provenance.}  The earlier 40--4--4 ``asymmetry'' was caused by assigning 18 rays shared by all three lineages to the first lineage by label order.  Exact projective accounting gives a symmetric decomposition: 21 common rays and 48 rays exclusive to each lineage; at the context level, 10 contexts are common and each lineage contributes 40 further contexts.  Each lineage therefore has 69 rays and 50 contexts, and the union has 165 rays and 130 contexts.  The revised Table~II marks all 21 common rays in black, and Table~III is left uncolored to avoid conflating shared and exclusive provenance.

\item \textbf{Completed the 20-ray orthogonality inventory.}  The concrete $D=4$ ray set admits 17 contexts: 10 imposed contexts, three automatic disjoint-support contexts, and four equal-modulus contexts.  Exact enumeration gives 16 separating two-valued states for the complete 20-ray hypergraph.  We also distinguish the nested 10-, 13-, and 17-context choices; the first two have the same 36 states, while the complete 17-context choice has 16.  The revised discussion explicitly respects the monotonicity that adding contexts cannot create new two-valued states.

\item \textbf{Repaired the mutual-reconstruction table.}  The gadget and Cabello sets share 15 rays, not 16.  The former $v_{1423}$/$v_{47}$ identification equated noncollinear vectors and has been replaced by two separate rank-three orthogonal-complement constructions.  Three additional reconstruction triples that were non-spanning or not orthogonal to their target have also been corrected.  The final table now displays five gadget-only and three Cabello-only rays, all reconstructible from the other set.

\item \textbf{Removed an incorrect $D=4$ MUB catalogue.}  The former appendix listed bases that did not satisfy the claimed pairwise overlap condition.  Because the forcing arguments require only the equal-modulus criterion for one vector relative to the computational basis, we removed that catalogue and retained a concise, verified qutrit MUB appendix.

\item \textbf{Style and figure captions.}  We shortened repetitive passages, corrected terminology and grammar, simplified the figure captions, and clarified that the drawings record incidence relations rather than establishing lattice properties by their planar layout.
\end{itemize}

We are grateful for the report: its recommendations prompted a clearer presentation and, through the accompanying audit, a materially more reliable manuscript.  We hope that the revised version addresses all of your concerns.

Sincerely,

Mirko Navara and Karl Svozil

\end{document}
